\section{当前选题面临的核心问题}
在已有系统的基础上，当前毕设如果继续按原有思路推荐，其主要问题不在于“系统功能不够丰富”，而在于选题主线尚未完成从应用系统到研究问题的提升。这使得当前选题在创新性、独立性和答辩的可解释性方面都存在明显风险。

\paragraph{1. 缺少一个足够明确的核心算法问题}

目前系统已经包含多Agent协作、显式状态机、HITL、\texttt{patch/replan}、蛋白质设计工作流等多个组成部分，但这些组成部分更多体现为一个较完整的执行平台，而不是一个被单独抽象和证明的算法对象。

换句话说，当前系统解决了“如何把流程跑起来”的问题，却没有清晰回答“系统依据什么原则做关键决策”这一更为核心的问题。尤其在以下方面，仍然缺少关键建模：

\begin{itemize}
	\item 如何刻画任务难度
	\item 如何估计执行风险
	\item 如何衡量工作量与代价
	\item 如何在不同风险和预算条件下调整执行策略
	\item 如何决定是否值得自动执行、请求人工接入或直接恢复/终止
\end{itemize}

虽然当前已有部分门控和打分机制，但更多属于启发式聚合和排序，而不是一个足够独立、可被论证为“核心算法”的研究对象。

因此，如果按照当前路径继续推进，很可能被评价为缺乏足够集中、足够有辨识度的算法贡献。

\paragraph{2. 当前选题仍然容易被归入“生命科学助手”子集}

已经有另一个同学的毕设是“生命科学助手”，同样采用多Agent架构。从架构相似性上来说，即便本系统已经具备一些机制，也很难天然摆脱“生命科学助手系统的一种更为复杂的实现”这一印象。

原因在于，当前系统的主要叙事仍然是：

$\texttt{用户提出生命科学任务} \rightarrow \texttt{多Agent规划与执行} \texttt{调工具} \rightarrow \texttt{返回结果}$

只要这个叙事不变，那么无论内部运行机制如何增强，它仍然有较大概率被看作生命科学助手的执行引擎、后端工作流或子模块。这对答辩会有很大风险。

因为该系统很难回答与已有的生命科学助手的本质区别、核心贡献是什么等这些问题。

\paragraph{3. 现有的增强机制更偏向工程健壮性, 而非算法创新}

当前系统中已经存在许多机制保证实验平台能力，但是从选题的角度去看，它们更多体现的是系统设计的正确性、运行治理等方面，没有进一步上升为统一的决策原理和算法问题。

\paragraph{4. 难度、工作量、风险尚未进入统一决策}

此前反复讨论的关键词包括难度、工作量、风险、不确定性、成本和人工介入时机等，这些词都指向一个共同问题：系统缺少的不是更多功能，而是一个能够解释并驱动关键治理动作的统一决策框架。

如果没有这一框架，那么就无法体现为什么某个任务更难、为什么某个任务需要人工介入、为什么某条恢复链路比另一条更合适等问题。

\paragraph{小结}

因此，当前选题面临的核心问题可以概括为三点：

第一，系统已经具备较完善的平台能力，但是尚未抽象出一个足够独立的核心算法问题。

第二，只要主叙事仍然停留在“生命科学领域的多Agent助手”，就难免被已经存在的父集系统覆盖。

第三，在当前时间约束下，必须尽快把研究对象从“继续扩展系统功能”转向“提炼一个可论证、可验证、可答辩”的核心问题。

也正因如此，后续的研究方向应该不再是“再加入什么能力”，而是究竟把当前系统上升为哪一个更为底层和独立的计算机专业课题的问题。
